%! Author = chtompki
%! Date = 1/14/26
% Example LaTeX document for GP111 - note % sign indicates a comment
\documentclass[12pt]{amsart}
% Default margins are too wide all the way around. I reset them here
\setlength{\hoffset}{-.75in}
\setlength{\textwidth}{6.5in}
\setlength{\voffset}{-.5in}
\setlength{\textheight}{9.0in}
\usepackage{amsmath}
\usepackage{amssymb}
\usepackage{amsthm}
\usepackage{pgfplots}
\usepackage{enumerate}
\usepackage{tikz}
\usetikzlibrary{shadows,arrows,arrows.meta,shapes,positioning,calc}% Required for the Circle and Triangle arrow tips
\usepackage{setspace}
\usepackage{siunitx}
\usepackage{enumitem}

\theoremstyle{conjecture}
\newtheorem{conjecture}{Conjecture}

\theoremstyle{definition}
\newtheorem{definition}{Definition}

\theoremstyle{question}
\newtheorem{question}{Question}

\theoremstyle{theorem}
\newtheorem{theorem}{Theorem}

\newenvironment{solution}
{\begin{proof}[Solution]}
{\end{proof}}

\title{Graph Theory:\\Homework 2\\ Book Section 1.2.}
\author{Rob Tompkins}
\date{\today}


\begin{document}

    \maketitle
    \noindent\textbf{1.} Suppose that every vertex of a simple graph $G$ has degree at least 3.
    Prove that $G$ has a cycle of even length.
    (Hint: Consider a maximal path $P=v_1 v_2 \cdots v_k$)

    \begin{proof}
        Let $P=v_1v_2\cdots v_k$ be a maximal path in $G$.  By maximality,
        every neighbor of $v_k$ lies on $P$; otherwise, a neighbor outside $P$
        could be appended to obtain a longer path.  Since $d(v_k)\geq 3$, the
        vertex $v_k$ has at least two neighbors on $P$ other than $v_{k-1}$.
        Choose two of them, say $v_i$ and $v_j$, with $i<j<k-1$.

        If either $k-i+1$ or $k-j+1$ is even, then the corresponding cycle
        \[
            v_i v_{i+1}\cdots v_kv_i
            \quad\text{or}\quad
            v_j v_{j+1}\cdots v_kv_j
        \]
        has even length.  Otherwise both of these numbers are odd, so $i$ and
        $j$ have the same parity.  In that case
        \[
            v_i v_{i+1}\cdots v_jv_kv_i
        \]
        is a cycle of length $(j-i)+2$, which is even because $j-i$ is even.
        Thus $G$ contains a cycle of even length.
    \end{proof}

    \noindent\textbf{2.} Prove that a bipartite graph has a unique bipartition
    (except for interchanging the two partite sets) if and only if it is connected.

    \begin{proof}
        First suppose that $G$ is connected, and let $(A,B)$ be a bipartition.
        Fix a vertex $r\in A$.  For every vertex $x$, any $r$--$x$ path
        alternates between $A$ and $B$.  Consequently, $x\in A$ exactly when
        $d(r,x)$ is even, and $x\in B$ exactly when $d(r,x)$ is odd.  Thus the
        choice of the part containing $r$ determines the entire bipartition.
        The only other choice interchanges $A$ and $B$, so the bipartition is
        unique up to interchange.

        Conversely, suppose that $G$ is disconnected and let $(A,B)$ be a
        bipartition.  Choose one connected component $H$ of $G$ and interchange
        the two partite sets only on $H$; that is, set
        \[
            A'=(A\setminus V(H))\cup(B\cap V(H)),\qquad
            B'=(B\setminus V(H))\cup(A\cap V(H)).
        \]
        There are no edges between $H$ and the other components, so $(A',B')$
        is also a bipartition.  Since at least one component was not switched,
        $(A',B')$ is neither $(A,B)$ nor $(B,A)$.  Hence a disconnected
        bipartite graph does not have a unique bipartition up to interchange.
    \end{proof}

    \noindent\textbf{3.} Let $v$ be a cut-vertex of a simple graph $G$.
    Prove that $\overline{G} - v$ is connected.

    \begin{proof}
        Because $v$ is a cut-vertex, $G-v$ has at least two connected
        components.  Let $x$ and $y$ be arbitrary vertices of $G-v$.  If they
        belong to different components of $G-v$, then $xy\notin E(G)$, and
        hence $xy\in E(\overline{G})$.  Thus they are adjacent in
        $\overline{G}-v$.


        If $x$ and $y$ belong to the same component of $G-v$, choose a vertex
        $z$ in a different component.  Neither $xz$ nor $yz$ is an edge of
        $G$, so both are edges of $\overline{G}$.  Therefore $xzy$ is a path
        in $\overline{G}-v$.  Every two vertices of $\overline{G}-v$ are thus
        joined by a path, so $\overline{G}-v$ is connected.
    \end{proof}

    \noindent\textbf{4.} Let $G$ be a simple graph having no isolated vertex and no induced
    subgraph with exactly two edges.
    Prove that $G$ is a complete graph.

    \begin{proof}
        We first show that $G$ is connected.  If it were disconnected, choose
        an edge $x_1x_2$ in one component and an edge $y_1y_2$ in another.
        Such edges exist because $G$ has no isolated vertices.  The subgraph
        induced by $\{x_1,x_2,y_1,y_2\}$ would have exactly the two edges
        $x_1x_2$ and $y_1y_2$, a contradiction.

        Now suppose that $G$ is not complete.  Choose two nonadjacent vertices
        $u$ and $w$.  Since $G$ is connected, there is a shortest
        $u$--$w$ path
        \[
            u=v_0,v_1,\ldots,v_m=w,
        \]
        where $m\geq 2$.  By the minimality of this path, $v_0v_2$ is not an
        edge.  Hence the subgraph induced by $\{v_0,v_1,v_2\}$ has exactly the
        two edges $v_0v_1$ and $v_1v_2$, again a contradiction.  Therefore
        $G$ must be complete.
    \end{proof}

    \noindent\textbf{5.} Let $G$ be a connected simple graph not having $P_4$, or $C_3$ as
    an induced sub-graph.
    Prove that $G$ is a biclique (complete bipartite graph).
    (Hint: Recall that the distance $d_G(u,v)$ between vertices $u$ and $v$ in $G$ is the shortest
    path between $u$ and $v$ [assuming that $G$ is connected].
    Fix an arbitrary vertex $v\in V(G)$.
    Partition $V(G)$ using distances to $v$ by letting $L_i=\{u\in V(G): d_g(v,u) = i\}$
    for $i\ge 0$.
    Why must $L_i=\null$ for $i\ge 3$.
    Why must $G$ be complete bipartite?)

    \begin{proof}
        Fix a vertex $v\in V(G)$ and, for each $i\geq 0$, let
        \[
            L_i=\{u\in V(G):d_G(v,u)=i\}.
        \]
        We first claim that $L_i=\varnothing$ for every $i\geq 3$.  Indeed, if
        $u\in L_i$ for some $i\geq 3$, then the first four vertices
        $v_0=v,v_1,v_2,v_3$ of a shortest $v$--$u$ path induce a $P_4$.
        There can be no edge joining two nonconsecutive vertices among them,
        because such an edge would give a shorter path from $v$ to $v_2$ or
        $v_3$.  This contradicts the hypothesis.  Therefore
        \[
            V(G)=L_0\cup L_1\cup L_2.
        \]

        The set $L_1$ is independent: an edge between two of its vertices,
        together with $v$, would induce a triangle.  The set $L_2$ is also
        independent.  To see this, suppose $x,y\in L_2$ and $xy$ is an edge.
        Choose $a\in L_1$ adjacent to $x$.  If $ay$ is an edge, then
        $a,x,y$ form a triangle.  If $ay$ is not an edge, then
        $v,a,x,y$ induce a $P_4$.  Both alternatives are impossible.

        Finally, every vertex of $L_1$ is adjacent to every vertex of $L_2$.
        Let $y\in L_2$, and choose $a\in L_1$ adjacent to $y$.  If some
        $b\in L_1$ were not adjacent to $y$, then $b,v,a,y$ would induce a
        $P_4$: the only possible additional edges are excluded because $L_1$
        is independent and because $y\notin L_1$.  This is impossible.

        Let $A=L_1$ and $B=L_0\cup L_2$.  Both $A$ and $B$ are independent,
        and every vertex of $A$ is adjacent to every vertex of $B$: vertices
        of $L_1$ are adjacent to $v$, and, as just proved, to every vertex of
        $L_2$.  Hence $G$ is the complete bipartite graph with partite sets
        $A$ and $B$.
    \end{proof}

\end{document}